\begin{proof}[Proof of \cref{obj:lem:radius-explicit-observed-kl}]
Fix \(t_0>0\).

1.  Apply \cref{obj:lem:observed-path-kl} with this value of \(t_0\).  It gives a constant \(K_0=K_0(t_0)>0\), and it remains only to supply the sequential-ignorability and stationary-start hypotheses for the signed-depth experiments over the range requested there.  For every auxiliary horizon \(T'\in\mathbb N\), terminal depth \(Q'\ge1\), and sign alternative \(v\in\{-1,+1\}\), the finite construction underlying \cref{obj:def:signed-depth-family} has the observed-state behavior kernel
\[
\Pr_v\!\left(A_t=a\mid X_{1:t},H_{1:t},A_{1:t-1},Y_{1:t-1}\right)=b(a\mid X_t),
\]
and the embedding into the observable-data experiment preserves this factorization, so \cref{obj:ass:sequential-ignorability} holds.  The same construction starts the joint state from its behavior stationary law,
\[
S_1\sim d_b,
\]
so \cref{obj:ass:stationary-start} holds for the same \((T',Q',v)\), including the zero-horizon case.  Hence \cref{obj:lem:observed-path-kl} supplies, for every \(\zeta>0\), \(C>1\), \(T\in\mathbb N\), and \(Q\ge1\), the sharpened observed-path bound
\[
\operatorname{KL}\!\left(
\mathbb P_{Q,+1}^{\mathrm{obs},T}\,\middle\|\,
\mathbb P_{Q,-1}^{\mathrm{obs},T}
\right)
\le
\frac{4c_0(t_0)^2K_0^2}{1-c_0(t_0)^2}\,
T\,\varepsilon_C^2\,\alpha(t_0)^{2Q}\,L(\zeta)^{-Q}.
\]
This KL estimate is one conjunct of the conclusion; the radius comparison is verified separately. % lean: CausalSmith.Stat.PomdpLatentOverlapMinimax.radius_explicit_observed_kl

2.  It remains to prove
\[
\frac{q_C}{2}\le \varepsilon_C\le q_C,
\qquad
q_C=\frac{C-1}{C},
\qquad
\varepsilon_C=\min\left\{\frac{C-1}{2},\frac12\right\}.
\]
Since \(C>1\), we have \(C>0\) and \(C-1\ge0\).  First suppose \(C\le2\).  Then
\[
\varepsilon_C=\frac{C-1}{2}.
\]
For the lower bound, \(q_C/2\le\varepsilon_C\) is
\[
\frac{C-1}{2C}\le\frac{C-1}{2},
\]
which follows from \(C>0\), \(C\ge1\), and \(C-1\ge0\).  For the upper bound, \(\varepsilon_C\le q_C\) is
\[
\frac{C-1}{2}\le\frac{C-1}{C}.
\]
Multiplying by \(2C>0\), this becomes
\[
C(C-1)\le2(C-1),
\]
and this follows from \(C-1\ge0\) and \(C\le2\). % lean: CausalSmith.Stat.PomdpLatentOverlapMinimax.radius_explicit_observed_kl

3.  Now suppose \(C>2\).  Then
\[
\varepsilon_C=\frac12.
\]
For the lower bound, \(q_C/2\le\varepsilon_C\) is
\[
\frac{C-1}{2C}\le\frac12,
\]
equivalently \(C-1\le C\), which is immediate.  For the upper bound, \(\varepsilon_C\le q_C\) is
\[
\frac12\le\frac{C-1}{C}.
\]
Multiplying by \(2C>0\), this becomes \(C\le2(C-1)\), equivalently \(2\le C\), which follows from \(C>2\). % lean: CausalSmith.Stat.PomdpLatentOverlapMinimax.radius_explicit_observed_kl

Combining the two cases gives the asserted radius comparison, and the KL bound is the one obtained in Step 1 with \(L(\zeta)=\exp(\zeta)\) and \(\alpha(t_0)=\exp(-1/t_0)\). % lean: CausalSmith.Stat.PomdpLatentOverlapMinimax.radius_explicit_observed_kl
\end{proof}
